\section{Zeit-Pulshöhen-Konverter}

Ein Zeit-Pulshöhen-Konverter (TAC) wandelt den Zeitunterschied zwischen einem Startsignal und einem Stopsignal in einen Puls mit einer zu der Zeitdifferenz proportionalen Spannung. Hierzu wird das Laden eines Kondensators mit konstantem Strom genutzt\cite[294]{leotechniques}.

Es wird ein \texttt{ORTEC} Modell 567 TAC verwendet. Dieser hat einen einstellbaren Auflösungszeitraum, welcher auf $\SI{100}{\nano\second}$ gesetzt ist\cite{tac}.
